% docs/manual/developer/chapters/05-registry.tex
% 第 5 章：registry 子系统

\chapter{registry 子系统}

\modname{data.registry} 是 \openbench{} 的 \emph{数据集与模型元信息} 中央索引。它把"哪个数据集叫什么、住哪、能算什么变量"从代码里搬出来，住进版本化的 YAML catalog，让用户级注册成为可能。本章详细解释 catalog 结构、Manager API，并给两个完整 walkthrough：注册新 reference 数据集 + 注册新 model profile。

\section{registry 鸟瞰}

\modname{data.registry} 包含：

\begin{itemize}
  \item \modname{schema}：catalog 字段的 dataclass（ReferenceDataset / ModelProfile / VariableMapping / FallbackVar）
  \item \modname{manager}：RegistryManager 公开 API + get\_registry() 单例
  \item \modname{scanner}：从 NC 文件自动扫元信息（years 范围、tim\_res、grid\_res）
  \item \modname{converter}：catalog 升级路径（旧格式 → 新格式）
  \item \modname{reference\_catalog.yaml}：~98 个 reference 数据集（系统级）
  \item \modname{model\_catalog.yaml}：21 个 model profile（系统级）
  \item \modname{reference\_profiles.yaml}：reference 数据集的额外 profile 信息
  \item \modname{station\_lists/}：站点列表 CSV 集合
\end{itemize}

\section{Catalog 加载优先级}

RegistryManager 加载 catalog 的顺序：

\begin{enumerate}
  \item 系统级（包内 \file{src/openbench/data/registry/*.yaml}）
  \item 用户级（platformdirs 的 user config 目录，如 \file{\~{}/Library/Application\ Support/openbench/} on macOS）
  \item 工作目录级（如果存在 \file{./openbench\_registry.yaml}）
\end{enumerate}

后加载的 \emph{覆盖} 先加载的同名 entry。这让用户能用 \cli{openbench data register} 注册私有数据集而不污染主仓库。

\section{ReferenceDataset schema}

\file{registry/schema.py} 定义 reference 入口的字段：

\begin{itemize}
  \item \yamlkey{name}：唯一标识（与 catalog key 一致）
  \item \yamlkey{description}：人类可读说明
  \item \yamlkey{category}：物理分类（Bio / Carbon / Water / Energy / Heat / Meteorology / Urban / Composite / Other）
  \item \yamlkey{data\_type}：\texttt{grid} 或 \texttt{stn}
  \item \yamlkey{root\_dir}：本地路径（支持 \texttt{\$ENV} 占位）
  \item \yamlkey{tim\_res}：\texttt{Month} / \texttt{Day} / \texttt{Hour} / \texttt{Year}
  \item \yamlkey{grid\_res}：浮点（仅 grid）
  \item \yamlkey{years}：\texttt{[start, end]}
  \item \yamlkey{timezone}：时区偏移
  \item \yamlkey{data\_groupby}：文件组织（\texttt{Year} / \texttt{Month} / \texttt{Single}）
  \item \yamlkey{variables}：dict，每个变量一个 VariableMapping
  \item \yamlkey{station\_matching}：dict（仅 stn），含 \texttt{area\_var} 等站点元字段
  \item \yamlkey{fulllist}：站点列表 CSV 路径（仅 stn）
  \item \yamlkey{\_provenance}：自动写入的元信息可信度标记
\end{itemize}

\section{VariableMapping 与 FallbackVar}

每个变量映射：

\begin{itemize}
  \item \yamlkey{varname}：NC 内变量名
  \item \yamlkey{varunit}：源单位
  \item \yamlkey{prefix} / \yamlkey{suffix}：文件名约定（多变量分文件时用）
  \item \yamlkey{sub\_dir}：相对于 root\_dir 的子目录
  \item \yamlkey{compute}：单位换算或层选择表达式
  \item \yamlkey{fallbacks}：FallbackVar 列表（同一物理量的多个候选 NC 名）
\end{itemize}

FallbackVar：

\begin{itemize}
  \item \yamlkey{varname}：fallback NC 变量名
  \item \yamlkey{varunit}：fallback 单位
  \item \yamlkey{convert}：fallback 的单位换算表达式（独立于主变量）
\end{itemize}

\section{ModelProfile schema}

类似 ReferenceDataset，但少几个字段（无 station\_matching、无 root\_dir，root\_dir 由用户在 simulation entry 写）：

\begin{itemize}
  \item \yamlkey{name}：模型名
  \item \yamlkey{data\_type}：\texttt{grid} / \texttt{stn}
  \item \yamlkey{tim\_res} / \yamlkey{grid\_res}
  \item \yamlkey{description}
  \item \yamlkey{variables}：标准变量名 → VariableMapping
  \item \yamlkey{time\_offset}：可选的时间戳偏移（处理"小时模型把月底标 -1day"等约定）
  \item \yamlkey{prefix\_fallback}：文件名匹配的多模式回退
\end{itemize}

\section{RegistryManager API}

主入口：

\begin{minted}{python}
from openbench.data.registry import RegistryManager
from openbench.data.registry.manager import get_registry

mgr = RegistryManager()    # 直接构造（默认加载所有 catalog）
mgr = get_registry()       # 单例（推荐，避免重复加载）
\end{minted}

\subsection{常用方法}

\begin{minted}{python}
# 列出全部
refs = mgr.list_references()       # list[ReferenceDataset]
models = mgr.list_models()          # list[ModelProfile]

# 按变量过滤
mgr.references_for_variable("Evapotranspiration")  # list

# 按名查
ref = mgr.get_reference("GLEAM_v4.2a_LowRes")  # ReferenceDataset | None
mp = mgr.get_model("CoLM2024")                  # ModelProfile | None

# 分辨率变体（base 名 → 变体字典）
variants = mgr.get_resolution_variants("GLEAM_v4.2a")
# {'LowRes': ReferenceDataset(...), 'MidRes': ReferenceDataset(...)}
\end{minted}

\subsection{可写 catalog}

\begin{minted}{python}
from openbench.data.registry.manager import (
    get_writable_reference_catalog_path,
    get_writable_model_catalog_path,
)

ref_path = get_writable_reference_catalog_path()
# ~/Library/Application Support/openbench/reference_catalog.yaml (macOS)
\end{minted}

\section{Scanner：自动元信息提取}

\file{registry/scanner.py} 在 \cli{openbench data register} 中被调用，从 NC 文件自动提取：

\begin{itemize}
  \item 时间分辨率（看 \texttt{time} 坐标的步长）
  \item 空间分辨率（看 \texttt{lat}/\texttt{lon} 步长）
  \item years 范围（time 坐标首尾）
  \item data\_type（含 \texttt{site} 或类似坐标 → stn；否则 grid）
  \item 可用变量列表
\end{itemize}

scanner 把推断结果写入 catalog 的 \yamlkey{\_provenance} 字段，后续 resolver 据此打 provenance 标记。

\section{Walkthrough A：添加新 Reference 数据集（10 步）}

假设你要把新 ET 数据集 \dataset{NewET\_v1} 注册到 \openbench{}：

\subsection{Step 1：准备数据}

\begin{minted}{bash}
# 假设 NC 文件结构
ls /data/myrefs/NewET_v1/
# NewET_v1_2010.nc  NewET_v1_2011.nc  ...
\end{minted}

每个 NC 文件含 \texttt{lat, lon, time, ET}（变量名 ET，单位 mm/day）。

\subsection{Step 2：选择标准变量名}

查 \openbench{} 的标准变量列表（用户卷附录 C 或 \cli{openbench data list}）。这里对应 \var{Evapotranspiration}。

\subsection{Step 3：决定 catalog 写入位置}

\begin{itemize}
  \item 私有用：\cli{openbench data register} 写到用户级 catalog
  \item 共享：手动编辑系统级 \file{src/openbench/data/registry/reference\_catalog.yaml} + 提 PR
\end{itemize}

这里走私有路径。

\subsection{Step 4：用 register 命令注册}

\begin{minted}{bash}
openbench data register NewET_v1 \
  --root-dir /data/myrefs/NewET_v1 \
  --data-type grid --grid-res 0.5 --tim-res Month \
  --years 2010 2020 --category Water \
  -v "Evapotranspiration:ET:mm/day"
\end{minted}

预期输出：

\begin{minted}{text}
✓ Created reference dataset 'NewET_v1' (1 variable)
Verify: openbench data list | grep NewET_v1
\end{minted}

\subsection{Step 5：验证 list 能看到}

\begin{minted}{bash}
openbench data list | grep NewET_v1
# NewET_v1   Water   grid   0.5°  2010-2020  1
\end{minted}

\subsection{Step 6：验证 path 正确}

\begin{minted}{bash}
openbench data path NewET_v1
# /data/myrefs/NewET_v1
\end{minted}

\subsection{Step 7：在 yaml 中引用}

\begin{minted}{yaml}
evaluation:
  variables:
    - Evapotranspiration

reference:
  Evapotranspiration: NewET_v1
\end{minted}

\subsection{Step 8：跑 check}

\begin{minted}{bash}
openbench check openbench.yaml
# ✓ Evapotranspiration → NewET_v1 (grid, Month, 0.5°)
\end{minted}

\subsection{Step 9：跑评估}

\begin{minted}{bash}
openbench run openbench.yaml --dry-run
\end{minted}

\subsection{Step 10：（可选）共享给社区}

如果数据集是公开的：

\begin{enumerate}
  \item Fork 主仓库
  \item 把 entry 添加到 \file{src/openbench/data/registry/reference\_catalog.yaml}（系统级）
  \item 加测试到 \file{tests/test\_registry/test\_manager.py}
  \item 提 PR
\end{enumerate}

\section{Walkthrough B：添加新 Model Profile（8 步）}

假设你要为新模型 \dataset{MyLandModel} 注册 profile：

\subsection{Step 1：列出模型输出的 NC 变量}

\begin{minted}{bash}
ncdump -h /data/mymodel/file.nc | grep "double\|float" | head -20
# float et_total ;
# float lh_flux ;
# float sh_flux ;
# float runoff ;
\end{minted}

\subsection{Step 2：建立标准变量 → NC 名 的映射}

\begin{longtable}{l l l}
  \toprule
  标准变量名 & NC 变量名 & 单位 \\
  \midrule
  \endhead
  \var{Evapotranspiration} & et\_total & kg/m2/s \\
  \var{Latent\_Heat} & lh\_flux & W/m2 \\
  \var{Sensible\_Heat} & sh\_flux & W/m2 \\
  \var{Total\_Runoff} & runoff & mm/s \\
  \bottomrule
\end{longtable}

\subsection{Step 3：决定单位换算}

\openbench{} 标准 ET 单位是 mm/day。NC 是 kg/m²/s。

水密度 1000 kg/m³，所以 1 kg/m²/s = 1 mm/s = 86400 mm/day。

\subsection{Step 4：用 register 命令注册}

\begin{minted}{bash}
openbench model register MyLandModel \
  --data-type grid --grid-res 0.5 --tim-res Month \
  --description "My new land surface model" \
  -v "Evapotranspiration:et_total:kg/m2/s" \
  -v "Latent_Heat:lh_flux:W/m2" \
  -v "Sensible_Heat:sh_flux:W/m2" \
  -v "Total_Runoff:runoff:mm/s"
\end{minted}

\subsection{Step 5：补 compute 表达式}

由于 register 命令不直接支持 compute 字段，要么手动编辑 catalog YAML，要么 register 后再 \cli{model show} 看缺什么再补。

手动编辑 \file{\~{}/Library/Application\ Support/openbench/model\_catalog.yaml}：

\begin{minted}{yaml}
MyLandModel:
  variables:
    Evapotranspiration:
      varname: et_total
      varunit: kg/m2/s
      compute: value * 86400   # → mm/day
    Total_Runoff:
      varname: runoff
      varunit: mm/s
      compute: value * 86400   # → mm/day
\end{minted}

\subsection{Step 6：验证 show}

\begin{minted}{bash}
openbench model show MyLandModel
\end{minted}

应能看到 4 个变量映射 + compute 表达式。

\subsection{Step 7：在 yaml 中引用}

\begin{minted}{yaml}
simulation:
  MyRun:
    model: MyLandModel
    root_dir: /data/mymodel
\end{minted}

\subsection{Step 8：跑评估验证}

\begin{minted}{bash}
openbench check openbench.yaml
openbench run openbench.yaml --dry-run
openbench run openbench.yaml
\end{minted}

如果输出数值不合理（ET 数量级离谱），检查 \texttt{compute} 表达式。

\section{Catalog YAML 完整示例}

完整 reference entry：

\begin{minted}{yaml}
GLEAM_v4.2a_LowRes:
  name: GLEAM_v4.2a_LowRes
  description: GLEAM v4.2a Low Resolution ET
  category: Water
  data_type: grid
  root_dir: ${OPENBENCH_REF_ROOT}/Grid/LowRes/Water/Evapotranspiration/GLEAM_v4.2a
  data_groupby: Year
  tim_res: Month
  grid_res: 0.5
  timezone: 0
  years: [1980, 2023]
  variables:
    Evapotranspiration:
      varname: E
      varunit: mm/day
      sub_dir: ""
      prefix: ""
      suffix: ""
    Transpiration:
      varname: Et
      varunit: mm/day
\end{minted}

完整 model entry：

\begin{minted}{yaml}
CoLM2024:
  name: CoLM2024
  description: Common Land Model 2024
  data_type: grid
  tim_res: Month
  grid_res: 0.5
  variables:
    Evapotranspiration:
      varname: f_lh_vap
      varunit: W/m2
      compute: value * 86400 / 2.5e6
      fallbacks:
        - varname: f_evapot
          varunit: mm/day
    Latent_Heat:
      varname: f_lh
      varunit: W/m2
\end{minted}

\section{下一步}

下一章详解 \modname{core} 评估引擎：metrics / scores / evaluation engines / comparison / 19 个 statistics 模块，含添加新 metric / score / statistic 的三个 walkthrough。
